\documentclass[11pt,letterpaper]{article}
\usepackage[T1]{fontenc}
\usepackage[utf8]{inputenc}
\usepackage{lmodern}
\usepackage[margin=0.93in,headheight=14pt]{geometry}
\usepackage{amsmath,amssymb,amsthm,mathtools}
\usepackage{microtype,booktabs,array}
\usepackage{enumitem}
\usepackage{fancyhdr}
\usepackage[hidelinks,bookmarksnumbered]{hyperref}
\usepackage{listings}
\hypersetup{pdftitle={Two Apery-transform supercongruences in OEIS A267220},
 pdfauthor={ChatGPT},pdfsubject={Proofs, parameter refinement, and the dyadic correction}}
\pagestyle{fancy}
\fancyhf{}
\fancyhead[L]{\small Ap\'ery-transform supercongruences}
\fancyhead[R]{\small OEIS A267220}
\fancyfoot[C]{\thepage}
\renewcommand{\headrulewidth}{0.3pt}
\setlength{\parskip}{0.25em}
\setlength{\emergencystretch}{2em}
\setlist{itemsep=0.2em,topsep=0.4em}
\numberwithin{equation}{section}
\newtheorem{theorem}{Theorem}[section]
\newtheorem{proposition}[theorem]{Proposition}
\newtheorem{lemma}[theorem]{Lemma}
\newtheorem{corollary}[theorem]{Corollary}
\theoremstyle{definition}
\newtheorem{definition}[theorem]{Definition}
\newtheorem{remark}[theorem]{Remark}
\newcommand{\Z}{\mathbb Z}
\newcommand{\Q}{\mathbb Q}
\newcommand{\Zp}{\Z_{(p)}}
\newcommand{\Ztwo}{\Z_{(2)}}
\newcommand{\dd}{\delta}
\newcommand{\CT}{\operatorname{CT}}
\newcommand{\Rev}{\operatorname{Rev}}
\newcommand{\Li}{\operatorname{Li}}
\newcommand{\ord}{\operatorname{ord}}
\newcommand{\vm}{\nu_p(m)}
\newcommand{\Bn}{\mathcal B}
\newcommand{\Wm}{\mathcal W}
\lstset{basicstyle=\small\ttfamily,columns=fullflexible,keepspaces=true,
 breaklines=true,frame=single,framerule=0.3pt,xleftmargin=0.4em,xrightmargin=0.4em}

\title{\textbf{Two Ap\'ery-transform supercongruences\\in OEIS A267220}\\[0.45em]
\large A coefficient proof, a parameter refinement,\\and the dyadic correction}
\author{Research note prepared for Vladimir Reshetnikov}
\date{September 20, 2026}

\begin{document}
\maketitle
\thispagestyle{empty}

\begin{abstract}
Two formulas contributed by Peter Bala to OEIS A267220 on October 17,
2024 assert second-order supercongruences for diagonal powers of the
exponential transform of the classical Ap\'ery numbers and of
its associated reversion transform. Both formulas are true for every
integer value of the parameter. We prove them by establishing the
second-order Gauss congruences for the Ap\'ery numbers directly from their
binomial sum, then giving a coefficient-level proof of preservation under
normalized framing. For odd primes the result gains an additional factor
of the prime dividing the parameter, and the restriction to primes at
least five can be removed. The second family satisfies the original
unsigned congruence also at the prime two; the first family requires an
explicit sign correction at the first dyadic descent. We prove the full
signed statement, treat zero and negative parameters, and give exact
counterexamples to two tempting stronger assertions. The general
framing-integrality mechanism is established prior work of Schwarz,
Vologodsky, and Walcher, not a novelty claim of this note. The purpose here
is to resolve the specific OEIS assertions with a complete, directly
checkable argument. A standard-library Python program and its exact
verification records accompany the article.
\end{abstract}

\section{The two conjectures and the result}\label{sec:target}

Let
\begin{equation}\label{eq:apery}
a_n=\sum_{k=0}^{n}\binom nk^2\binom{n+k}{k}^2\qquad(n\geq0)
\end{equation}
be the classical Ap\'ery numbers, OEIS A005259 \cite{oeis-apery}.
Their initial values are
\[
1,\ 5,\ 73,\ 1445,\ 33001,\ 819005,\ 21460825,\ldots.
\]
The two series relevant to the problem are
\begin{equation}\label{eq:A-F}
A(x)=\exp\!\left(\sum_{n\geq1}\frac{a_n}{n}x^n\right),
\qquad
F(x)=\frac{x}{\Rev(xA(x))}.
\end{equation}
Here $\Rev$ denotes compositional inversion at zero. Since
$xA(x)=x+O(x^2)$, this inverse exists uniquely over $\Q[[x]]$.
The quotient defining $F$ is a well-defined formal power series with
constant coefficient one. Integrality of both series will be proved
below; it is not assumed in the argument.

The coefficients of $A$ form OEIS A267220. The entry includes two
conjectures, both explicitly attributed to Peter Bala and dated
October 17, 2024 \cite{oeis-target}. In our notation they concern
\begin{equation}\label{eq:u-v}
u_m(n)=[x^n]A(x)^{mn},\qquad
v_m(n)=[x^n]F(x)^{mn}\qquad(m\in\Z,\ n\geq0),
\end{equation}
and assert
\begin{equation}\label{eq:oeis-conjectures}
\begin{split}
u_m(np^r)&\equiv u_m(np^{r-1})\pmod{p^{2r}},\\
v_m(np^r)&\equiv v_m(np^{r-1})\pmod{p^{2r}}
\end{split}
\end{equation}
for every prime $p\geq5$ and all positive integers $n,r$.
These assertions were still labeled conjectural when the entry was
checked on September 20, 2026.

The notation $a_n$ in this article always denotes the Ap\'ery numbers,
\emph{not} the coefficients of A267220. In particular, these are two
related but different sequences. Also,
\[
u_m(0)=v_m(0)=1
\]
for every $m$, whereas $u_0(n)=v_0(n)=0$ for $n\geq1$.
This distinction matters when normalizing by the parameter.

\begin{theorem}[Resolution and refinements]\label{thm:main}
Both congruences \eqref{eq:oeis-conjectures} hold for every integer $m$.
More precisely, let $N\geq1$, let $p\mid N$ be prime, and put
$R=\nu_p(N)\geq1$.
\begin{enumerate}[label=\textup{(\alph*)}]
\item If $p$ is odd and $m\neq0$, then
\begin{equation}\label{eq:main-odd}
\begin{split}
\nu_p\bigl(u_m(N)-u_m(N/p)\bigr)&\geq2R+\nu_p(m),\\
\nu_p\bigl(v_m(N)-v_m(N/p)\bigr)&\geq2R+\nu_p(m).
\end{split}
\end{equation}
\item The unsigned congruence for the second family holds at every prime,
including $2$:
\begin{equation}\label{eq:main-v-all}
v_m(N)\equiv v_m(N/p)\pmod{p^{2R}}.
\end{equation}
\item For every prime $p$ and $m\neq0$, the signed refinements are
\begin{equation}\label{eq:main-signed}
\begin{aligned}
(-1)^{mN}u_m(N)
&\equiv(-1)^{mN/p}u_m(N/p)
&&\pmod{p^{2R+\nu_p(m)}},\\
(-1)^{(m-1)N}v_m(N)
&\equiv(-1)^{(m-1)N/p}v_m(N/p)
&&\pmod{p^{2R+\nu_p(m)}}.
\end{aligned}
\end{equation}
\end{enumerate}
For $m=0$, all the positive-index terms in these statements vanish.
\end{theorem}

We use $\nu_p(0)=+\infty$, and $\nu_p(m)=\nu_p(|m|)$ for nonzero
negative integers. Substituting $N=np^r$ gives
$R=r+\nu_p(n)$, so part (a) proves a modulus at least as strong as the
one in OEIS even when $p$ divides $n$. The theorem does not tacitly
assume $p\nmid n$ or $p\nmid m$.

\subsection*{What is established here, and what is prior work?}
The nonlinear preservation principle behind the proof is the
\emph{integrality of framing} for $2$-functions. It is proved in general
by Schwarz, Vologodsky, and Walcher \cite{svw2013,svw2017}; their 2017
paper includes a local coefficient proof. Our proof uses that known
mechanism in an explicit normalization adapted to the two OEIS families.
No claim is made that the general preservation theorem, or its connection
to dilogarithms, is new. Nor does a conjectural label in OEIS establish
that no earlier specialist has deduced this particular application.

The contribution of this note is a self-contained resolution of the
stated formulas, with all parameter values and the prime two handled
explicitly, together with reproducible computations. The proof does not
appeal to a computer search as a substitute for an all-index argument.
Its dependency chain is
\[
\begin{gathered}
\text{Ap\'ery binomial sum}\ \Longrightarrow\
\text{second-order Gauss congruences}\ \Longrightarrow\
\text{integral }A,\\
\text{normalized coefficient transform}\ \Longrightarrow\
\text{framing congruences}\ \Longrightarrow\
\text{both OEIS assertions}.
\end{gathered}
\]

\section{Formal arithmetic and the dilogarithm criterion}\label{sec:arithmetic}

Everything below takes place in formal power series rings. No analytic
convergence or choice of a complex logarithm is involved. Set
\[
\dd=x\frac{d}{dx},\qquad
\CT\!\left(\sum_{j\geq j_0}q_jx^j\right)=q_0.
\]
The constant-term functional is applied only to Laurent series bounded
below. It satisfies
\begin{equation}\label{eq:ct-parts}
\CT(\dd G)=0,\qquad
\CT(G\dd J)=-\CT(J\dd G).
\end{equation}
Each relevant constant term depends on only finitely many coefficients.

For a prime $p$, write
\[
\Zp=\{a/b\in\Q:p\nmid b\}.
\]
A congruence between elements of $\Zp$ has its usual valuation meaning;
a congruence between series means the congruence of every coefficient.
Working in this localization is sufficient: completeness is unnecessary.

\begin{definition}\label{def:gauss}
An integer sequence $(\alpha_n)_{n\geq1}$ satisfies the
\emph{second-order Gauss congruences} if
\begin{equation}\label{eq:gauss-general}
\alpha_N\equiv\alpha_{N/p}\pmod{p^{2\nu_p(N)}}
\end{equation}
for every prime $p$ and every positive $N$ divisible by $p$.
\end{definition}

\begin{lemma}[M\"obius and dilogarithm criterion]\label{lem:mobius}
For an integer sequence $(\alpha_n)_{n\geq1}$, the following are
equivalent:
\begin{enumerate}[label=\textup{(\roman*)}]
\item The sequence satisfies \eqref{eq:gauss-general}.
\item The numbers
\begin{equation}\label{eq:c-mobius}
c_n=\frac1{n^2}\sum_{d\mid n}\mu(n/d)\alpha_d
\end{equation}
are integers for every $n\geq1$.
\item There is an expansion with integer coefficients
\begin{equation}\label{eq:dilog-expansion}
\sum_{n\geq1}\frac{\alpha_n}{n^2}x^n
=\sum_{d\geq1}c_d\Li_2(x^d),
\qquad \Li_2(z)=\sum_{j\geq1}\frac{z^j}{j^2}.
\end{equation}
\end{enumerate}
The $c_d$ in \textup{(iii)} are uniquely given by \eqref{eq:c-mobius}.
\end{lemma}

\begin{proof}
Fix $n=p^rh$ with $p\nmid h$. In the numerator of
\eqref{eq:c-mobius}, the only nonzero terms have divisor $p$-adic
exponent $r$ or $r-1$. Pairing them gives
\[
\sum_{d\mid n}\mu(n/d)\alpha_d
=\sum_{e\mid h}\mu(h/e)
 \bigl(\alpha_{p^re}-\alpha_{p^{r-1}e}\bigr).
\]
Under (i), this numerator is divisible by $p^{2r}$. Doing this for every
prime divisor of $n$ proves divisibility by $n^2$, hence (ii).

M\"obius inversion gives
\begin{equation}\label{eq:divisor-recovery}
\alpha_n=\sum_{d\mid n}d^2c_d.
\end{equation}
If every $c_d$ is an integer, subtracting the formula for
$\alpha_{N/p}$ from that for $\alpha_N$ leaves only divisors $d$ with
$\nu_p(d)=\nu_p(N)$. Every surviving term contains
$p^{2\nu_p(N)}$, proving (i). Finally, coefficient extraction from the
right side of \eqref{eq:dilog-expansion} gives
$n^{-2}\sum_{d\mid n}d^2c_d$, proving the equivalence with (iii).
\end{proof}

One immediate benefit of this stronger integrality criterion is an Euler
product. Define
\[
W(x)=\sum_{n\geq1}\frac{\alpha_n}{n^2}x^n,
\qquad H(x)=\dd W(x)=\sum_{n\geq1}\frac{\alpha_n}{n}x^n.
\]
Taking $\dd$ of \eqref{eq:dilog-expansion} and then exponentiating gives
\begin{equation}\label{eq:euler}
\exp H(x)=\prod_{d\geq1}(1-x^d)^{-d c_d}
\in1+x\Z[[x]].
\end{equation}
This identity is formal: to compute the coefficient of $x^n$, only
factors with $d\leq n$ are needed. The exponents may be negative or
positive; any integer power of a unit in $1+x\Z[[x]]$ is integral.

A second formulation will be central to the coefficient proof.

\begin{lemma}[Integral Frobenius defect]\label{lem:defect}
If $(\alpha_n)$ satisfies the second-order Gauss congruences, then for
every prime $p$ the series
\begin{equation}\label{eq:defect}
E_p(x)=W(x)-\frac1{p^2}W(x^p)
\end{equation}
belongs to $x\Zp[[x]]$.
\end{lemma}

\begin{proof}
If $p\nmid n$, the coefficient is $\alpha_n/n^2\in\Zp$. If $p\mid n$,
it is $(\alpha_n-\alpha_{n/p})/n^2$, again in $\Zp$ by the hypothesis.
\end{proof}

\section{The required Ap\'ery congruence from the binomial sum}\label{sec:apery}

We now return to the sequence $a_n$ in \eqref{eq:apery}. Only a
second-order congruence is needed. It admits a short proof that works
uniformly at all primes, including $2$ and $3$.

\begin{proposition}\label{prop:apery-gauss}
For every prime $p$ and positive integer $N$ divisible by $p$,
\begin{equation}\label{eq:apery-gauss}
a_N\equiv a_{N/p}\pmod{p^{2\nu_p(N)}}.
\end{equation}
\end{proposition}

\begin{proof}
Put $R=\nu_p(N)$ and
\[
P(N,k)=\binom Nk\binom{N+k}{k}.
\]
When $p\nmid k$, the identity
\[
\binom Nk=\frac Nk\binom{N-1}{k-1}
\]
shows that $\nu_p(P(N,k))\geq R$. Such a summand $P(N,k)^2$ therefore
vanishes modulo $p^{2R}$.

It remains to compare the summands with $k=p\ell$. For $k\geq1$, an
exact product identity is
\begin{equation}\label{eq:P-product}
P(N,k)=(-1)^{k-1}\frac Nk\left(1+\frac Nk\right)
  \prod_{j=1}^{k-1}\left(1-\frac{N^2}{j^2}\right).
\end{equation}
To see this, write the binomial product as
\[
\frac{N(N+k)}{k^2}
\prod_{j=1}^{k-1}\frac{(N-j)(N+j)}{j^2}
\]
and extract the signs. Separating the factors with $p\mid j$ in
\eqref{eq:P-product} yields
\begin{equation}\label{eq:P-split}
P(N,p\ell)=(-1)^{(p-1)\ell}P(N/p,\ell)
 \prod_{\substack{1\leq j<p\ell\\p\nmid j}}
 \left(1-\frac{N^2}{j^2}\right).
\end{equation}
Each factor in the remaining product belongs to
$1+p^{2R}\Zp$. Squaring removes the sign even when $p=2$, so
\[
P(N,p\ell)^2\equiv P(N/p,\ell)^2\pmod{p^{2R}}.
\]
The $k=0$ summand is one on both sides. Summing over $0\leq k\leq N$
now gives exactly \eqref{eq:apery-gauss}.
\end{proof}

Consequently, Lemma~\ref{lem:mobius} applies to $a_n$, and the series
$A$ in \eqref{eq:A-F} is integral by \eqref{eq:euler}. Its primitive
dilogarithm coefficients begin
\begin{equation}\label{eq:initial-c}
c_1=5,\qquad c_2=17,\qquad c_3=160,\qquad
c_4=2058,\qquad c_5=32760.
\end{equation}
For example, $c_4=(a_4-a_2)/16=2058$. Thus
\[
A(x)=(1-x)^{-5}(1-x^2)^{-34}(1-x^3)^{-480}
       (1-x^4)^{-8232}\cdots.
\]
The arithmetic input to the rest of the proof is now complete. We have
not used a stronger Ap\'ery supercongruence as a black box.

\section{An integral normalization of the two transforms}\label{sec:normalization}

From now on set
\begin{equation}\label{eq:WHD}
W=\sum_{n\geq1}\frac{a_n}{n^2}x^n,\qquad
H=\dd W=\log A,\qquad
D=\dd H=\sum_{n\geq1}a_nx^n.
\end{equation}
The series $D$ has constant coefficient zero; in particular, it is the
Ap\'ery generating series with its initial term $a_0=1$ removed.

For every $m\in\Z$, define $X_m(y)$ by
\begin{equation}\label{eq:X}
X_m(y)=yA(X_m(y))^m,
\qquad\text{equivalently}\qquad y=xA(x)^{-m}.
\end{equation}
The inverse exists uniquely in $y+y^2\Z[[y]]$: a series
$x+O(x^2)$ with integer coefficients has an integer compositional inverse,
as is seen by solving successively for its coefficients. This includes
negative $m$, since $A$ and every integer power of $A$ are integral units.

Define the normalized transform by
\begin{equation}\label{eq:B-definition}
\Bn_m(y)=\frac{D(X_m(y))}{1-mD(X_m(y))}
        =\sum_{n\geq1}b_m(n)y^n.
\end{equation}
The denominator is a unit of constant coefficient one. Therefore
\begin{equation}\label{eq:B-integer}
b_m(n)\in\Z\quad(m\in\Z,\ n\geq1),
\qquad b_0(n)=a_n.
\end{equation}
This integral definition is preferable to starting with a quotient by
$m$. In particular, it remains valid when a prime under consideration
divides $m$.

\begin{proposition}[Coefficient normalization]\label{prop:B-coefficient}
For $m\neq0$ and $n\geq1$,
\begin{equation}\label{eq:B-coefficient}
b_m(n)=\frac1m[x^n]A(x)^{mn}.
\end{equation}
Consequently, $u_m(n)=m b_m(n)$ for every integer $m$ and every $n\geq1$.
\end{proposition}

\begin{proof}
Logarithmic differentiation of $y=xA(x)^{-m}$ gives
\[
\frac{d\log y}{d\log x}=1-mD(x),
\]
so the chain rule implies
\begin{equation}\label{eq:chain}
\dd_yH(X_m(y))=\Bn_m(y).
\end{equation}
The formal Lagrange inversion formula (proved in
Appendix~\ref{app:lagrange}) gives
\[
[y^n]H(X_m(y))
=\frac1n[x^{n-1}]H'(x)A(x)^{mn}.
\]
For $m\neq0$, use $(A^{mn})'=mn H'A^{mn}$. The last display becomes
\[
[y^n]H(X_m(y))=\frac1{mn}[x^n]A(x)^{mn}.
\]
Multiplication by $n$, as in \eqref{eq:chain}, proves the assertion.
At $m=0$, the identity $u_0(n)=0=m b_0(n)$ holds separately.
\end{proof}

The relation between the two conjectures is particularly simple after
this normalization.

\begin{proposition}[The reversion family is a shifted transform]\label{prop:V}
The series $F$ is integral and satisfies
\begin{equation}\label{eq:F-X}
F(y)=A(X_{-1}(y)).
\end{equation}
For every integer $m$ and every $n\geq1$,
\begin{equation}\label{eq:V-B}
v_m(n)=m b_{m-1}(n).
\end{equation}
In particular, $v_1(n)=a_n$.
\end{proposition}

\begin{proof}
Equation \eqref{eq:X} with $m=-1$ says
$X_{-1}(y)A(X_{-1}(y))=y$, hence
$X_{-1}=\Rev(xA(x))$. Dividing by $X_{-1}$ proves
\eqref{eq:F-X}, including its integrality.

Apply Lagrange inversion to $A(X_{-1}(y))^{mn}$. For $n\geq1$ it gives
\begin{equation}\label{eq:v-lagrange}
[y^n]A(X_{-1}(y))^{mn}
=m[x^{n-1}]A'(x)A(x)^{(m-1)n-1}.
\end{equation}
If $m\neq1$, differentiating $A^{(m-1)n}$ turns this into
\[
\frac{m}{m-1}[x^n]A(x)^{(m-1)n}=m b_{m-1}(n).
\]
For $m=1$, the right side of \eqref{eq:v-lagrange} is
$[x^{n-1}]A'/A=[x^{n-1}]H'=a_n=b_0(n)$. Thus the apparent singularity
at $m=1$ is removable, and the asserted formula holds for all $m$.
\end{proof}

\begin{remark}[Why division modulo a prime is not enough]
The formula $v_m(n)=\frac{m}{m-1}u_{m-1}(n)$ for $m\neq1$ is an
identity over $\Q$, but it cannot safely be used to transfer a congruence
by dividing by $m-1$ when $p\mid m-1$. Similarly, dividing a congruence
for $u_m$ by $m$ may lose powers of $p$. The proof below establishes the
congruence for the integer $b_m$ itself. This is why no exceptional
parameter classes need to be excluded.
\end{remark}

\section{The normalized framing congruence at odd primes}\label{sec:odd}

The next theorem is the coefficient form of the established
framing-integrality principle. Its proof is included rather than invoked
from the general theory \cite{svw2013,svw2017}.

\begin{theorem}[Normalized framing congruences]\label{thm:normalized}
Let $b_m$ be as in \eqref{eq:B-definition}. For every integer $m$, define
\begin{equation}\label{eq:beta}
\beta_m(n)=(-1)^{mn}b_m(n)\qquad(n\geq1).
\end{equation}
Then $\beta_m$ satisfies the second-order Gauss congruences at every
prime. Equivalently, for $p\mid N$ and $R=\nu_p(N)$,
\begin{equation}\label{eq:beta-gauss}
(-1)^{mN}b_m(N)\equiv(-1)^{mN/p}b_m(N/p)\pmod{p^{2R}}.
\end{equation}
For odd $p$ this is the ordinary unsigned congruence for $b_m$.
At $p=2$ the unsigned congruence also holds whenever $m$ is even or
$R\geq2$.
\end{theorem}

We prove the odd-prime assertion in this section and the remaining cases
in Section~\ref{sec:two}. These arguments use only the integer
coefficients $a_n$ and their second-order Gauss congruences. Thus the same
proof establishes the theorem for any input sequence satisfying
Definition~\ref{def:gauss}.

\subsection{An exact constant-term expression}

Fix $p\mid N$, put $R=\nu_p(N)$, and first take $m\neq0$. For brevity
write $E=E_p$. Lemma~\ref{lem:defect} gives
\[
E=W(x)-p^{-2}W(x^p)\in x\Zp[[x]].
\]
Differentiating and using $H=\dd W$ yields
\begin{equation}\label{eq:frob-A}
\dd E=H(x)-p^{-1}H(x^p),\qquad
A(x^p)=A(x)^p\exp(-p\dd E).
\end{equation}
Introduce the integral Laurent series
\begin{equation}\label{eq:K}
K(x)=x^{-N}A(x)^{mN}\in x^{-N}\Z[[x]].
\end{equation}
Since substitution $x\mapsto x^p$ turns the coefficient of $x^{N/p}$
into the coefficient of $x^N$, Proposition~\ref{prop:B-coefficient}
and \eqref{eq:frob-A} give the exact identity
\begin{equation}\label{eq:master}
b_m(N)-b_m(N/p)
=\frac1m\CT\!\left(K\left(1-\exp(-mN\dd E)\right)\right).
\end{equation}
Expanding the exponential gives
\begin{equation}\label{eq:master-expanded}
b_m(N)-b_m(N/p)
=\sum_{k\geq1}\frac{(-1)^{k+1}m^{k-1}N^k}{k!}
       \CT\bigl(K(\dd E)^k\bigr).
\end{equation}
Only $k\leq N$ can contribute, because $\dd E$ has no constant term.
Thus no question about convergence of the infinite displayed sum arises.
The factor $m$ has disappeared from the denominator of each summand:
this is the crucial normalization when $p\mid m$.

\subsection{The linear term gains a second power of the index}

The $k=1$ term seems initially to contain only $N$, not $N^2$.
Constant-term integration by parts supplies the missing factor. Namely,
\begin{equation}\label{eq:delta-K}
\dd K=N(mD-1)K,
\end{equation}
and therefore
\begin{equation}\label{eq:linear}
\begin{aligned}
N\CT(K\dd E)
&=-N\CT(E\dd K)\\
&=N^2\CT\bigl(K(1-mD)E\bigr)
\in p^{2R}\Zp.
\end{aligned}
\end{equation}
All coefficients inside the final constant term are $p$-integral.
This identity, rather than a coarse exponential estimate alone, is what
produces a second-order congruence.

\subsection{The nonlinear terms at odd primes}

For an odd prime and $k\geq2$,
\begin{equation}\label{eq:factorial-odd}
\nu_p(k!)\leq k-2.
\end{equation}
For $k=2$ this is immediate. For $k\geq3$, Legendre's formula gives
$\nu_p(k!)\leq(k-1)/(p-1)\leq k-2$.
Since $\nu_p(m)\geq0$ and $R\geq1$,
\[
\begin{aligned}
\nu_p\left(\frac{m^{k-1}N^k}{k!}\right)
&\geq kR-(k-2)\\
&=2R+(k-2)(R-1)\geq2R.
\end{aligned}
\]
Also $\CT(K(\dd E)^k)\in\Zp$. Every nonlinear term in
\eqref{eq:master-expanded} is therefore divisible by $p^{2R}$.
Together with \eqref{eq:linear}, this proves
\begin{equation}\label{eq:b-odd}
b_m(N)\equiv b_m(N/p)\pmod{p^{2R}}
\qquad(p\text{ odd},\ m\neq0).
\end{equation}
For $m=0$, this is precisely Proposition~\ref{prop:apery-gauss}.
Finally, $N$ and $N/p$ have the same parity when $p$ is odd, so the two
signs in \eqref{eq:beta-gauss} agree. This proves the odd-prime part of
Theorem~\ref{thm:normalized}.

\section{The prime two and the exact sign correction}\label{sec:two}

The odd-prime factorial bound fails at $p=2$. Replacing it indiscriminately
by the dyadic bound would lose information and would not prove the
asserted modulus. There are three distinct cases. Throughout this section
$E=E_2$, $K=x^{-N}A^{mN}$, and $R=\nu_2(N)\geq1$.
The case $m=0$ is already covered by the input Ap\'ery congruence.

\subsection{Even parameters}

Suppose $m\neq0$ is even. For $k\geq2$, the bound
$\nu_2(k!)\leq k-1$ gives
\[
\nu_2\left(\frac{m^{k-1}N^k}{k!}\right)
\geq(k-1)+kR-(k-1)=kR\geq2R.
\]
The linear term is handled by \eqref{eq:linear}. Hence
\begin{equation}\label{eq:two-even}
b_m(N)\equiv b_m(N/2)\pmod{2^{2R}}
\qquad(m\text{ even}).
\end{equation}
The signs in \eqref{eq:beta-gauss} are both positive in this case.

\subsection{Odd parameters beyond the first dyadic descent}

Suppose $m$ is odd and $R\geq2$. For $k\geq3$,
\[
\nu_2\left(\frac{m^{k-1}N^k}{k!}\right)
\geq kR-(k-1)\geq2R.
\]
The last inequality follows from
$(k-2)R\geq k-1$, valid for $R\geq2$ and $k\geq3$.
The only problematic nonlinear term is $k=2$, whose scalar coefficient
$mN^2/2$ has valuation $2R-1$.

Write $E=\sum_{j\geq1}e_jx^j$, with $e_j\in\Ztwo$.
Reduction modulo two gives
\begin{equation}\label{eq:delta-E-square}
(\dd E)^2\equiv\sum_{\substack{j\geq1\\j\text{ odd}}}
 e_j^2x^{2j}\pmod2.
\end{equation}
For odd $j$, applying $\CT\dd$ to $Kx^{2j}$ and using
\eqref{eq:delta-K} yields
\begin{equation}\label{eq:ct-even}
2j\,\CT(Kx^{2j})
=N\CT\bigl(K(1-mD)x^{2j}\bigr).
\end{equation}
Since $j$ is odd, the right side shows
$\CT(Kx^{2j})\in2^{R-1}\Ztwo\subseteq2\Ztwo$.
Equation \eqref{eq:delta-E-square} consequently implies
\[
\CT\bigl(K(\dd E)^2\bigr)\in2\Ztwo.
\]
This supplies the missing factor two in the quadratic term. The linear
term again follows from \eqref{eq:linear}, proving
\begin{equation}\label{eq:two-higher}
b_m(N)\equiv b_m(N/2)\pmod{2^{2R}}
\qquad(m\text{ odd},\ R\geq2).
\end{equation}
Here $N$ and $N/2$ are both even, so this is also the signed congruence.

\subsection{Odd parameters at the first dyadic descent}

It remains to take $m$ odd and $N=2h$ with $h$ odd. The correct relation
is now a \emph{sum}:
\begin{equation}\label{eq:two-plus-goal}
b_m(2h)+b_m(h)\equiv0\pmod4.
\end{equation}
Using the same derivation as in \eqref{eq:master}, but adding instead of
subtracting, we obtain
\begin{equation}\label{eq:plus-master}
b_m(N)+b_m(h)
=\frac1m\CT\!\left(K\left(1+\exp(-mN\dd E)\right)\right).
\end{equation}
Let $s_2(k)$ denote the number of ones in the binary expansion of $k$.
Legendre's formula at two is
\[
\nu_2(k!)=k-s_2(k).
\]
As $mN$ has valuation one, the coefficient $( -mN)^k/k!$ has valuation
$s_2(k)$. Modulo four, only positive exponents $k$ that are powers of two
survive. Each surviving scalar is congruent to two modulo four, since it
has valuation exactly one. Thus, coefficientwise in $\Ztwo[[x]]$,
\begin{equation}\label{eq:binary-exponential}
1+\exp(-mN\dd E)
\equiv2\left(1+\sum_{s\geq0}(\dd E)^{2^s}\right)\pmod4.
\end{equation}
The sum on the right is locally finite in degree.

We next identify its reduction modulo two. In $\dd E$, terms of even
degree vanish modulo two. For odd $j$, $e_j=a_j/j^2$, and hence
\[
\dd E\equiv\sum_{j\text{ odd}}a_jx^j\pmod2.
\]
Frobenius in characteristic two gives
\[
\sum_{s\geq0}(\dd E)^{2^s}
\equiv\sum_{s\geq0}\sum_{j\text{ odd}}a_jx^{j2^s}\pmod2.
\]
The input congruences imply $a_{j2^s}\equiv a_j\pmod2$.
Every positive integer is uniquely $j2^s$ with $j$ odd. Therefore
\begin{equation}\label{eq:binary-D}
\sum_{s\geq0}(\dd E)^{2^s}\equiv D(x)\pmod2.
\end{equation}
Combining \eqref{eq:plus-master}--\eqref{eq:binary-D} gives
\[
b_m(N)+b_m(h)\equiv\frac2m\CT\bigl(K(1+D)\bigr)\pmod4.
\]
Because $m$ is odd,
$1+D\equiv1-mD\pmod2$. But \eqref{eq:delta-K} and $\CT\dd K=0$
give the exact identity
\begin{equation}\label{eq:ct-zero}
\CT\bigl(K(1-mD)\bigr)=0.
\end{equation}
It follows that $\CT(K(1+D))$ is even, proving
\eqref{eq:two-plus-goal}. Since $(-1)^{mN}=1$ and
$(-1)^{mh}=-1$, this is exactly the missing case of
\eqref{eq:beta-gauss}. Theorem~\ref{thm:normalized} is proved.

\section{Deduction of both OEIS conjectures}\label{sec:deduction}

We now prove Theorem~\ref{thm:main} and make clear where the two families
behave differently at two.

\subsection{Odd primes and the parameter factor}

For $p$ odd, Theorem~\ref{thm:normalized} and the identities
\[
u_m(n)=m b_m(n),\qquad v_m(n)=m b_{m-1}(n)
\]
immediately imply
\[
\begin{split}
u_m(N)-u_m(N/p)&\in m p^{2R}\Zp,\\
v_m(N)-v_m(N/p)&\in m p^{2R}\Zp.
\end{split}
\]
This is \eqref{eq:main-odd}. In particular, the originally stated range
$p\geq5$ may be replaced by all odd primes.

Multiplying the all-prime signed congruence for $b_m$ by $m$ proves the
first part of \eqref{eq:main-signed}. Applying the same congruence to
$b_{m-1}$ and multiplying by $m$ proves its second part. Notice that the
sign parameter for $v_m$ is $m-1$, but the extra valuation comes from
$m$, not from $m-1$.

\subsection{Why the second family works unsigned at two}

When $m$ is odd, the normalized parameter $m-1$ is even. Thus
\eqref{eq:two-even} immediately proves
$v_m(N)\equiv v_m(N/2)\pmod{2^{2R}}$.

When $m$ is even and $R\geq2$, \eqref{eq:two-higher} gives the same
result, in fact with the extra factor $2^{\nu_2(m)}$.
It remains to consider a nonzero even $m$ and $N=2h$ with $h$ odd.
Now $m-1$ is odd, so the signed theorem says
\[
b_{m-1}(2h)+b_{m-1}(h)\in4\Z.
\]
Consequently,
\begin{equation}\label{eq:V2-boundary}
\begin{aligned}
v_m(2h)-v_m(h)
&=m\bigl(b_{m-1}(2h)+b_{m-1}(h)\bigr)
  -2m b_{m-1}(h)\\
&\in2^{\nu_2(m)+1}\Z\subseteq4\Z.
\end{aligned}
\end{equation}
This proves \eqref{eq:main-v-all} in the last case. The zero parameter
has already been separated, so Theorem~\ref{thm:main} is complete.

\subsection{A precise dyadic summary}

For the first family, the unsigned parameter-refined bound
\[
\nu_2\bigl(u_m(N)-u_m(N/2)\bigr)\geq2R+\nu_2(m)
\]
holds for nonzero even $m$, and also for odd $m$ when $R\geq2$.
When $m$ is odd and $R=1$, replace the difference by a sum.

For the second family, the same parameter-refined bound holds whenever
$m$ is odd or $R\geq2$. At the remaining first-descent case with $m$
even, \eqref{eq:V2-boundary} gives the always-valid unsigned exponent
$\nu_2(m)+1$, while the signed statement retains exponent
$2+\nu_2(m)$. The ordinary exponent two holds without exception.
The examples in Section~\ref{sec:examples} show why this distinction
cannot simply be suppressed.

\section{Framing as a change of coordinate}\label{sec:framing}

The preceding proof can be understood without geometric background.
Nevertheless, its structure becomes more transparent when the normalized
transforms are organized as a family of coordinate changes.

\subsection{The additive law for the transformations}

For an integral unit $A$, set
\[
\mathfrak T_m(A)(y)=A(X_m(y))=:A_m(y),
\qquad y=xA(x)^{-m}.
\]
If a second transformation has parameter $\ell$, its coordinate is
\[
z=yA_m(y)^{-\ell}=xA(x)^{-(m+\ell)}.
\]
Uniqueness of formal inversion therefore proves
\begin{equation}\label{eq:group}
\mathfrak T_\ell\bigl(\mathfrak T_m(A)\bigr)
=\mathfrak T_{m+\ell}(A),\qquad
\mathfrak T_0(A)=A.
\end{equation}
Thus $F=\mathfrak T_{-1}(A)$ is not an unrelated construction. It is
one step in the same integer-parameter family. The shift from $b_m$ to
$b_{m-1}$ in \eqref{eq:V-B} reflects precisely this additive law.

\subsection{The transformed potential}

Define
\begin{equation}\label{eq:framed-potential}
\Wm_m(y)=W(X_m(y))-\frac m2 H(X_m(y))^2.
\end{equation}
Using $d\log y=(1-mD(x))d\log x$, differentiation of the right side
gives
\[
\dd_y\Wm_m(y)=H(X_m(y)),\qquad
\dd_y^2\Wm_m(y)=\Bn_m(y).
\]
All three series have zero constant term where appropriate, so
\begin{equation}\label{eq:potential-coeff}
\Wm_m(y)=\sum_{n\geq1}\frac{b_m(n)}{n^2}y^n.
\end{equation}
The signed congruence theorem and Lemma~\ref{lem:mobius} now imply an
integral dilogarithm expansion
\begin{equation}\label{eq:framed-dilog}
\Wm_m((-1)^m y)=\sum_{d\geq1}c_{m,d}\Li_2(y^d),
\end{equation}
where
\begin{equation}\label{eq:framed-c}
c_{m,d}=\frac1{d^2}\sum_{e\mid d}\mu(d/e)(-1)^{me}b_m(e)\in\Z.
\end{equation}
In particular,
\begin{equation}\label{eq:framed-Euler}
A_m((-1)^m y)=\prod_{d\geq1}(1-y^d)^{-d c_{m,d}}.
\end{equation}
The sign in this coordinate is invisible to odd-prime descent, but it
is essential for the all-prime dilogarithm expansion.

This is the framing-integrality interpretation of the coefficient
proof. Its general mathematical theory and geometric setting are the
subject of \cite{svw2013,svw2017}. Higher-order congruences under more
special framing hypotheses are studied in \cite{mueller}; no such
additional hypothesis is invoked here. In particular, the arithmetic of
the input must not be assumed to transfer at every stronger order merely
because second order transfers.

\section{Examples and limits of stronger assertions}\label{sec:examples}

\subsection{Low-degree values and exceptional parameters}

Direct expansion gives
\begin{align*}
A(x)&=1+5x+49x^2+685x^3+11807x^4+232771x^5+O(x^6),\\
F(x)&=1+5x+24x^2+200x^3+2430x^4+36096x^5+O(x^6).
\end{align*}
These values agree with the target entry \cite{oeis-target}.
The normalized coefficients are already nontrivial polynomials in the
parameter:
\begin{equation}\label{eq:small-polynomials}
\begin{split}
b_m(1)&=5,\\
b_m(2)&=73+50m,\\
b_m(3)&=\frac{2890+3285m+1125m^2}{2}.
\end{split}
\end{equation}
The third expression has rational polynomial coefficients but takes
integer values at every integer $m$, as guaranteed structurally by
\eqref{eq:B-definition}. One should not confuse an integer-valued
polynomial with a polynomial all of whose coefficients are integers.

\begin{center}
\small
\begin{tabular}{r r r r r r}
\toprule
$m$ & $b_m(1)$ & $b_m(2)$ & $b_m(3)$ & $b_m(4)$ & $b_m(5)$\\
\midrule
$-2$ & 5 & $-27$ & 410 & $-1915$ & 53130\\
$-1$ & 5 & 23 & 365 & 6343 & 123255\\
0 & 5 & 73 & 1445 & 33001 & 819005\\
1 & 5 & 123 & 3650 & 118059 & 4015380\\
2 & 5 & 173 & 6980 & 301517 & 13540505\\
3 & 5 & 223 & 11435 & 623375 & 35175630\\
\bottomrule
\end{tabular}
\end{center}

The cases $m=0$ and $m=1$ now have concrete interpretations.
The normalized family at $m=0$ is the Ap\'ery sequence, but the original
$u_0$ and $v_0$ vanish at positive indices. The reversion identity
$v_1(n)=a_n$ is an exact coefficient identity, not just a congruence.
Furthermore, $v_2(n)=2u_1(n)$ for every positive $n$.

For comparison, some primitive coefficients from
\eqref{eq:framed-c} are
\begin{center}
\small
\begin{tabular}{r r r r r}
\toprule
$m$ & $c_{m,1}$ & $c_{m,2}$ & $c_{m,3}$ & $c_{m,4}$\\
\midrule
$-1$ & $-5$ & 7 & $-40$ & 395\\
0 & 5 & 17 & 160 & 2058\\
1 & $-5$ & 32 & $-405$ & 7371\\
2 & 5 & 42 & 775 & 18834\\
\bottomrule
\end{tabular}
\end{center}
Negative coefficients present no difficulty: integrality, rather than
positivity, is the property under consideration.

\subsection{The first unsigned family cannot include two unconditionally}

At $m=1$, $N=2$,
\[
u_1(2)-u_1(1)=123-5=118\not\equiv0\pmod4.
\]
Thus a blanket unsigned extension of the first OEIS assertion to $p=2$
is false. The signed statement instead gives
\[
u_1(2)+u_1(1)=128\equiv0\pmod4,
\]
exactly as predicted by the first-descent proof.

The stronger parameter factor at two also has a genuine boundary for
the second family:
\[
v_2(2)-v_2(1)=246-10=236\equiv0\pmod4,
\qquad236\not\equiv0\pmod8.
\]
Since $2R+\nu_2(m)=3$ in this example, the unsigned parameter-refined
bound cannot hold uniformly at two. The original second-order
congruence still holds, and the signed sum is $256$, divisible by $8$.

\subsection{Second order cannot uniformly be replaced by third order}

A small odd-prime example also limits the possible strengthening:
\begin{equation}\label{eq:sharpness}
\begin{aligned}
u_1(7)-u_1(1)
&=5082313271\\
&=7^2\cdot103720679,
\qquad103720679\equiv6\pmod7.
\end{aligned}
\end{equation}
The valuation is exactly two, not three. Therefore a universal
replacement of $p^{2r}$ by $p^{3r}$ in the first conjecture is false.
Because $v_2(n)=2u_1(n)$ and $7\nmid2$, the same example rules out that
replacement for the second family as well. This is a counterexample to
a uniform higher-order assertion, not a claim that no individual
parameters or primes can exhibit stronger divisibility.

\section{Exact computation and reproducibility}\label{sec:computation}

The companion file \texttt{verify.py} uses only the Python standard
library and arbitrary-precision integers. Every division that is
mathematically supposed to be exact is checked with a remainder test.
The general proof above is independent of the computations; the program
checks identities, edge cases, and representative congruences by routes
that expose indexing or normalization mistakes.

\subsection{Coefficient algorithms}

For a fixed integer $q$, write $A(x)^q=\sum_{j\geq0}h_jx^j$.
Logarithmic differentiation gives
\begin{equation}\label{eq:coefficient-recurrence}
h_0=1,\qquad
h_j=\frac qj\sum_{k=1}^j a_k h_{j-k}\quad(j\geq1).
\end{equation}
The quotient is an integer because $A^q$ is an integral unit. This
recurrence is valid for positive, zero, and negative $q$.
To obtain $u_m(N)$, use the \emph{fixed} exponent $q=mN$ throughout the
computation through degree $N$. The exponent must not be changed to
$mj$ as the internal coefficient index $j$ varies.

For $m\neq0$, compute $b_m(N)$ by exact division of $h_N$ by $m$;
for $m=0$, use $a_N$. Then compute $v_m(N)=m b_{m-1}(N)$.
This requires $O(N^2)$ integer arithmetic operations and $O(N)$ stored
integers for one coefficient at a fixed parameter and target index.
These are arithmetic-operation bounds, not unit-cost bit-complexity
claims for the large integers involved.

The input numbers are generated directly from the binomial definition
using
\[
P(n,0)=1,\qquad
P(n,k)=P(n,k-1)\frac{(n-k+1)(n+k)}{k^2},
\qquad a_n=\sum_{k=0}^nP(n,k)^2.
\]
They are independently compared against the recurrence recorded in
A005259 \cite{oeis-apery}:
\begin{equation}\label{eq:apery-recurrence}
(n+1)^3a_{n+1}
=(34n^3+51n^2+27n+5)a_n-n^3a_{n-1},
\qquad a_0=1,\ a_1=5.
\end{equation}
This familiar recurrence is used only as a cross-check; neither its
validity nor a proof of it is needed for the congruence argument.

For independent low-degree reversion checks, Lagrange inversion gives
\[
[y^n]F(y)=
\begin{cases}
5,&n=1,\\[2pt]
\dfrac1{1-n}[x^n]A(x)^{1-n},&n\geq2.
\end{cases}
\]
The program verifies $F(y)=A(y/F(y))$, both directions of compositional
inversion, and direct polynomial powers of $A$ and $F$ through degree
12. These direct power checks do not assume the formulas $u_m=m b_m$
and $v_m=m b_{m-1}$ as their coefficient-computation method.

\subsection{The included run}

The archived run used maximum index $250$, original parameters
$-6\leq m\leq6$, and normalized parameters $-7\leq m\leq6$.
For the transformed congruences it used primes
$2,3,5,7,11,13$ and indices $N=hp^r\leq250$ with
$1\leq h\leq6$ and $r\geq1$. Duplicate prime--index pairs were removed,
leaving 71 pairs. This is a specified finite sample, not every possible
transformed parameter and index up to 250.

The input Ap\'ery congruences were checked more exhaustively: every prime
divisor of every positive $N\leq250$ was used. The completed run reported
\textbf{8,142 successful checks}. The grouped counts are:
\begin{center}
\small
\begin{tabular}{p{0.78\linewidth}r}
\toprule
Check category & Count\\
\midrule
Ap\'ery definition versus three-term recurrence & 251\\
Input Ap\'ery second-order congruences & 473\\
Input primitive dilogarithm coefficients & 250\\
Functional equation, two inverse identities, Euler product & 4\\
Direct coefficient identities for $u_m$ and $v_m$ & 312\\
Signed normalized congruences at all selected primes & 994\\
Unsigned normalized congruences in the applicable cases & 973\\
Unsigned $v_m$ congruences, including the prime two & 923\\
Unsigned $u_m$ congruences for odd primes or even $m$ & 815\\
Signed parameter-refined congruences for both families & 1704\\
Odd-prime parameter-refined congruences for both families & 1272\\
Signed transformed primitive dilogarithm coefficients & 168\\
Boundary and higher-order counterexamples & 3\\
\midrule
Total & 8142\\
\bottomrule
\end{tabular}
\end{center}

The JSON report gives all 71 prime--index pairs, the ungrouped category
counts, the exact counterexample values, and the arithmetic environment.
The CSV file records sample values of the sequences. No floating-point
approximation or probabilistic primality test is used.

\subsection{Reproduction commands}

From the archive directory, with Python 3.9 or later, run
\begin{lstlisting}
python verify.py --max-index 250 --output-dir verification
\end{lstlisting}
A larger finite test can be requested, for example, with
\begin{lstlisting}
python verify.py --max-index 500 --parameter-radius 8 \
    --multipliers 8 --primes 2 3 5 7 11 13 17 \
    --output-dir verification-500
\end{lstlisting}
Only the default run described above is represented by the archived
verification results. To rebuild the article, use
\begin{lstlisting}
latexmk -pdf -interaction=nonstopmode -halt-on-error article.tex
\end{lstlisting}
The source is also compatible with two ordinary \texttt{pdflatex} runs.
The computational artifact is not a proof-assistant formalization; the
proof is the mathematical argument in this article.

\section{Conclusion}

The two parameterized assertions in A267220 follow from the same
normalized transform, with a shift of one in the framing parameter.
The central arithmetic fact is preservation of second-order Gauss
congruences after this normalization. Its coefficient proof isolates
the only linear term that appears to lose one power of the index and
recovers that power by formal integration by parts.

For the stated OEIS conjectures the conclusion is unconditional: every
integer parameter, every allowed prime, and every positive pair of
indices is covered. In fact, all odd primes work with an additional
parameter valuation, and the second family works unsigned at two.
The dyadic sign correction is exact, not an artifact of an insufficient
estimate, as the smallest counterexample demonstrates. Finally,
second order is genuinely the strongest order that can be asserted
uniformly for these families without additional restrictions.

\appendix
\section{A formal Lagrange inversion proof}\label{app:lagrange}

For completeness, the coefficient formula used above is established
here in the precise formal setting required.

\begin{lemma}\label{lem:lagrange}
Let $\phi(x)\in1+x\Q[[x]]$ and let $X(y)$ be the unique series satisfying
$X=y\phi(X)$. For any $G(x)\in\Q[[x]]$ and $n\geq1$,
\begin{equation}\label{eq:lagrange}
[y^n]G(X(y))=\frac1n[x^{n-1}]G'(x)\phi(x)^n.
\end{equation}
\end{lemma}

\begin{proof}
Let $\operatorname{Res}_x$ select the coefficient of $x^{-1}$ from a
formal differential. Residues are unchanged by a substitution
$y=x+O(x^2)$. This can be checked on monomials: derivatives have zero
residue, and $y'/y=x^{-1}$ plus a power series.

Substitute $y=x/\phi(x)$ in the coefficient residue. Then
\[
\begin{aligned}
[y^n]G(X(y))
&=\operatorname{Res}_y G(X(y))y^{-n-1}\,dy\\
&=\operatorname{Res}_x
G(x)x^{-n-1}\phi(x)^n
\left(1-\frac{x\phi'(x)}{\phi(x)}\right)\,dx.
\end{aligned}
\]
The residue of $d(G(x)x^{-n}\phi(x)^n)$ is zero. Expanding that
derivative and rearranging shows that the last display equals
\[
\frac1n\operatorname{Res}_x G'(x)x^{-n}\phi(x)^n\,dx
=\frac1n[x^{n-1}]G'(x)\phi(x)^n,
\]
as required.
\end{proof}

\begin{thebibliography}{9}
\bibitem{oeis-target}
OEIS Foundation Inc., \emph{A267220: Expansion of
$\exp(\sum_{n\geq1}\mathrm{A005259}(n)x^n/n)$}.
The two supercongruence conjectures are attributed to Peter Bala,
October 17, 2024. Entry checked September 20, 2026.
\url{https://oeis.org/A267220}.

\bibitem{oeis-apery}
OEIS Foundation Inc., \emph{A005259: Ap\'ery numbers}.
Entry checked September 20, 2026.
\url{https://oeis.org/A005259}.

\bibitem{svw2013}
A.~Schwarz, V.~Vologodsky, and J.~Walcher,
\emph{Framing the Di-Logarithm (over $\Z$)},
arXiv:1306.4298 (2013).
\url{https://arxiv.org/abs/1306.4298}.

\bibitem{svw2017}
A.~Schwarz, V.~Vologodsky, and J.~Walcher,
\emph{Integrality of Framing and Geometric Origin of 2-functions
(with algebraic coefficients)},
arXiv:1702.07135, version 2 (2017), especially Section 3.
\url{https://arxiv.org/abs/1702.07135}.

\bibitem{mueller}
L.~Felipe M\"uller,
\emph{Wolstenholme Type Congruences and Framing of Rational 2-Functions},
arXiv:2104.10754 (2021).
\url{https://arxiv.org/abs/2104.10754}.
\end{thebibliography}

\end{document}
